\uuid{122a}
\titre{ {TITRE} }
\theme{Polynôme et racines}
\auteur{Quentin LIARD}
\datecreate{2025-02-13}
\organisation{AMSCC}

\contenu{

\texte{ On pose le polynôme $P$ définie par $$P(X)=X^4+6X^3+13X^2+12X+4$$}
 

 \begin{enumerate}
	\item \question{Vérifier que $1$ est racine de multiplicité deux.}
	\indication{}
    \reponse{Le reste de la division euclidienne de $P$ par $X^2+2X+1=(X+1)^2$ est nul.}
    \item \question{Déterminer l'autre racine réelle de multiplicité deux.}
    	\indication{}
    \reponse{$P(X)=(X+1)^2(X^2+4X+4)=(X+1)^2(X+2)^2$ et ainsi $-2$ est racine double.}
\end{enumerate}

}